\providecommand{\gardetitre}{\par\vskip 0pt plus 9\baselineskip\penalty-300\vskip 0pt plus -9\baselineskip}
\chapter{Modèles}
\label{ann:modeles}

Ces modèles sont ceux de la mission ANPP. Ils se recopient et s'adaptent. Aucun ne dépasse une page une fois
rempli, sauf la note de cadrage et le schéma directeur lui-même.

\gardetitre
\section*{B1. Note de cadrage (Chapitre~\ref{chap:cadrer})}

\begin{tableaucours}{cL{4.0cm}L{7.4cm}}{%
    \tableauHeader{N°} &
    \tableauHeader{Section} &
    \tableauHeader{Contenu}
}
1 & Fait déclencheur & Trois lignes~: ce qui s'est passé, à quelle date \\
2 & Commanditaire et approbation & Qui commande, quel organe approuve, par quel acte \\
3 & Périmètre & Entités, métiers, objets. Liste des exclusions datées et motivées \\
4 & Horizon et ambition & Années couvertes. Remettre à niveau, aligner ou transformer \\
5 & Instances & Composition nominative, présidence, rythme \\
6 & Parties prenantes & Grille pouvoir-intérêt, conduite pour chaque groupe \\
7 & Responsabilités & Matrice par livrable, un seul approbateur \\
8 & Calendrier et livrables & Phases, dates, comités, plafond de pages \\
9 & Communication & À qui, quoi, quand, par quel moyen \\
10 & Signature & Commanditaire, date \\
\end{tableaucours}

\gardetitre
\section*{B2. Fiche d'application (Chapitre~\ref{chap:diagnostic})}

\begin{tableaucours}{L{4.4cm}L{7.6cm}}{%
    \tableauHeader{Rubrique} &
    \tableauHeader{Exemple~: liquidation des droits}
}
Nom et métier servi & LIQUID, direction des prestations \\
Responsable côté métier & Directeur des prestations \\
Utilisateurs & 64 agents, siège et antennes \\
Éditeur ou auteur & Développement interne \\
Mise en service & 2004 \\
Technologie & Langage et base qui ne sont plus maintenus \\
Maintenance & Deux agents de la division \\
Coût annuel & 38~millions, salaires compris \\
Données personnelles & Oui~: identité, carrière, droits \\
Notes & Valeur 4,5. Santé 1,5. Décision~: refondre \\
\end{tableaucours}

\gardetitre
\section*{B3. Fiche de projet (Chapitre~\ref{chap:portefeuille})}

\begin{tableaucours}{L{4.4cm}L{7.6cm}}{%
    \tableauHeader{Rubrique} &
    \tableauHeader{Exemple~: P07, référentiel des bénéficiaires}
}
Résultat & Un fichier unique des 186~000 bénéficiaires, maître pour toutes les applications \\
Objectif servi & Ne payer que les ayants droit \\
Responsable & Directeur des prestations \\
Investissement & 340~millions (60 en 2027, 280 en 2028) \\
Coût annuel ensuite & 48~millions à partir de 2029 \\
Durée & 20~mois, passation comprise \\
Dépendances & Aucune en amont. Les trois lots de la refonte en dépendent \\
Ce qu'il remplace & Onze fichiers de bénéficiaires \\
Risque principal & Reprise des données défaillante. Parade~: fiabilisation préalable \\
Sécurité, continuité, données & Analyse de risque, reprise sur le site de secours, formalités de la loi n°~2024/017 \\
Jalon d'arrêt & Fin de la fiabilisation~: moins de 1~\% de doublons, sinon report \\
\end{tableaucours}

\gardetitre
\section*{B4. Fiche d'indicateur (Chapitre~\ref{chap:piloter})}

\begin{tableaucours}{L{4.4cm}L{7.6cm}}{%
    \tableauHeader{Rubrique} &
    \tableauHeader{Exemple~: paiements à date}
}
Formule & Prestations payées au plus tard à la date prévue, divisées par les prestations dues \\
Source & Journal des paiements \\
Responsable de la mesure & Cellule de portefeuille \\
Fréquence & Chaque mois \\
Mesure de départ & 91~\% en 2026 \\
Cible & 99~\% en 2031 \\
\end{tableaucours}

\gardetitre
\section*{B5. Sommaire type d'un schéma directeur}

Soixante pages au plus, hors annexes. Une synthèse de quatre pages l'accompagne pour l'organe qui approuve.

\begin{tableaucours}{cL{5.0cm}cL{4.6cm}}{%
    \tableauHeader{N°} &
    \tableauHeader{Partie} &
    \tableauHeader{Pages} &
    \tableauHeader{Chapitre du cours}
}
1 & Synthèse pour décision & 4 & Tous \\
2 & Mandat, périmètre, gouvernance de la démarche & 4 & \ref{chap:cadrer} \\
3 & Stratégie de l'organisation et cadre de référence & 5 & \ref{chap:cameroun} et \ref{chap:droit} \\
4 & Diagnostic~: dix constats & 8 & \ref{chap:diagnostic} \\
5 & Cible~: capacités, principes, plan fonctionnel, scénario retenu & 10 & \ref{chap:cible} \\
6 & Portefeuille et trajectoire & 10 & \ref{chap:portefeuille} \\
7 & Financement, achats, dotation & 8 & \ref{chap:financer} \\
8 & Pilotage et révision & 5 & \ref{chap:piloter} \\
9 & Risques, sécurité, continuité, conformité & 6 & \ref{chap:maitriser} \\
 & Annexes~: fiches, inventaire, registres, projets écartés &  &  \\
\end{tableaucours}

\gardetitre
\section*{B6. Page de décision pour un comité}

Chaque comité reçoit trois jours avant une page par décision. Elle contient cinq lignes~: la décision demandée,
les faits et leurs pièces, les options avec leur coût, la recommandation, ce qui se passe si l'on ne décide pas.
